
\documentclass{article}
\usepackage{colortbl}
\usepackage{makecell}
\usepackage{multirow}
\usepackage{supertabular}

\begin{document}

\newcounter{utterance}

\centering \large Interaction Transcript for game `hot\_air\_balloon', experiment `air\_balloon\_survival\_en\_reasoning off\_hard', episode 3 with llama{-}3.3{-}70b{-}instruct{-}t1.0.
\vspace{24pt}

{ \footnotesize  \setcounter{utterance}{1}
\setlength{\tabcolsep}{0pt}
\begin{supertabular}{c@{$\;$}|p{.15\linewidth}@{}p{.15\linewidth}p{.15\linewidth}p{.15\linewidth}p{.15\linewidth}p{.15\linewidth}}
    \# & $\;$A & \multicolumn{4}{c}{Game Master} & $\;\:$B\\
    \hline

    \theutterance \stepcounter{utterance}  
    & & \multicolumn{4}{p{0.6\linewidth}}{
        \cellcolor[rgb]{0.9,0.9,0.9}{
            \makecell[{{p{\linewidth}}}]{
                \texttt{\tiny{[P1$\langle$GM]}}
                \texttt{Stai partecipando a un gioco di negoziazione collaborativa.} \\
\\ 
\texttt{Insieme a un altro partecipante, dovete accordarvi su un unico set di oggetti da conservare. Ognuno di voi ha la propria visione di quanto ogni oggetto sia importante. Non sai come l'altro partecipante valuta gli oggetti. Inoltre, ti viene fornito lo sforzo che ogni oggetto richiede.} \\
\texttt{Potete accordarvi su un set solo se lo sforzo totale degli oggetti selezionati non supera un limite condiviso:} \\
\\ 
\texttt{LIMITE = 3629} \\
\\ 
\texttt{Ecco i valori di sforzo dei singoli oggetti:} \\
\\ 
\texttt{Sforzo oggetto = \{"A21": 5, "B93": 793, "A11": 21, "B33": 842, "C08": 313, "B39": 842, "B96": 631, "C85": 226, "A82": 87, "B31": 761, "B20": 230, "A60": 287, "B09": 698, "C81": 641, "B21": 882\}} \\
\\ 
\texttt{Ecco la tua visione personale sull'importanza di ogni oggetto:} \\
\\ 
\texttt{Valori di importanza dell'oggetto = \{"A21": 138, "B93": 583, "A11": 868, "B33": 822, "C08": 783, "B39": 65, "B96": 262, "C85": 121, "A82": 508, "B31": 780, "B20": 461, "A60": 484, "B09": 668, "C81": 389, "B21": 808\}} \\
\\ 
\texttt{Obiettivo:} \\
\\ 
\texttt{Il tuo obiettivo è negoziare un set condiviso di oggetti che ti avvantaggi il più possibile (cioè, massimizzi l'importanza totale per TE), rimanendo entro il LIMITEE. Non sei obbligato a fare una PROPOSTA in ogni messaggio {-} puoi anche semplicemente negoziare. Tutte le tattiche sono permesse!} \\
\\ 
\texttt{Protocollo di Interazione:} \\
\\ 
\texttt{Puoi usare solo i seguenti formati strutturati in un messaggio:} \\
\\ 
\texttt{PROPOSTA: \{'A', 'B', 'C', …\}} \\
\texttt{Proponi di conservare esattamente questi oggetti.} \\
\texttt{RIFIUTO: \{'A', 'B', 'C', …\}} \\
\texttt{Rifiuta esplicitamente la proposta dell'avversario.} \\
\texttt{ARGOMENTAZIONE: \{'...'\}} \\
\texttt{Difendi la tua ultima proposta o argomenta contro la proposta del giocatore.} \\
\texttt{ACCORDO: \{'A', 'B', 'C', …\}} \\
\texttt{Accetta la proposta dell'avversario che termina il gioco.} \\
\\ 
\\ 
\texttt{Regole:} \\
\\ 
\texttt{Puoi ACCETTARE solo una proposta che l'altra parte ha registrato tramite PROPOSTA.} \\
\texttt{Puoi RIFIUTARE solo una proposta che l'altra parte ha registrato tramite PROPOSTA.} \\
\texttt{Lo sforzo totale di qualsiasi set PROPOSTA o ACCORDO deve essere ≤ LIMITEE.} \\
\texttt{NON rivelare i tuoi punteggi di importanza nascosti.} \\
\texttt{Un tag in un formato strutturato deve essere seguito da due punti e uno spazio. L'argomento deve essere un set python contenente 0 o più stringhe.} \\
\texttt{Quindi, deve essere della forma TAG: \{...\}} \\
\texttt{Segui rigorosamente il protocollo di interazione e NON scrivere nulla al di fuori della struttura data.} \\
\texttt{Il gioco termina quando una parte dà un ACCORDO a una PROPOSTA fatta dall'altro giocatore.} \\
\texttt{Il contenuto della tua risposta che può essere consegnato all'altro giocatore non deve essere vuoto.} \\
\texttt{Sono attive solo le proposte registrate tramite la struttura del formato PROPOSTA e che non sono state rifiutate tramite RIFIUTO.} \\
\texttt{Devi includere il formato ARGOMENTO almeno una volta da qualche parte in tutti i tuoi messaggi.} \\
            }
        }
    }
    & \\ \\

    \theutterance \stepcounter{utterance}  
    & \multicolumn{4}{p{0.6\linewidth}}{
        \cellcolor[rgb]{0.9,0.9,0.9}{
            \makecell[{{p{\linewidth}}}]{
                \texttt{\tiny{[P1$\rangle$GM]}}
                \texttt{ARGOMENTAZIONE: \{'A21', 'A11', 'C08', 'A82', 'B20'\} } \\
\texttt{Sulla base dei valori di sforzo forniti, questi oggetti sembrano avere un buon equilibrio tra sforzo e potenziale importanza, considerando la loro varietà e copertura di categorie diverse. Includono oggetti con basso sforzo come 'A21' e 'A11', che potrebbero essere facilmente aggiunti senza superare il limite, e oggetti con sforzo più alto come 'C08' e 'B20', che potrebbero essere cruciali per raggiungere un alto livello di importanza complessiva. 'A82' aggiunge varietà senza un onere eccessivo in termini di sforzo. Questa selezione sembra offrire un buon punto di partenza per future negoziazioni, bilanciando l'importanza relativa con il limite di sforzo totale.} \\
            }
        }
    }
    & & \\ \\

    \theutterance \stepcounter{utterance}  
    & & \multicolumn{4}{p{0.6\linewidth}}{
        \cellcolor[rgb]{0.9,0.9,0.9}{
            \makecell[{{p{\linewidth}}}]{
                \texttt{\tiny{[P1$\langle$GM]}}
                \texttt{La tua risposta conteneva una sequenza non taggata o hai usato RAGIONAMENTO STRATEGICO più di una volta. Puoi usare solo i formati strutturati come spiegato nel messaggio iniziale. Devono essere tutti della forma TAG: \{...\}} \\
            }
        }
    }
    & \\ \\

    \theutterance \stepcounter{utterance}  
    & \multicolumn{4}{p{0.6\linewidth}}{
        \cellcolor[rgb]{0.9,0.9,0.9}{
            \makecell[{{p{\linewidth}}}]{
                \texttt{\tiny{[P1$\rangle$GM]}}
                \texttt{ARGOMENTAZIONE: \{'A21', 'A11', 'C08'\} } \\
\texttt{PROPOSTA: \{'A21', 'A11', 'C08'\}} \\
            }
        }
    }
    & & \\ \\

    \theutterance \stepcounter{utterance}  
    & & \multicolumn{4}{p{0.6\linewidth}}{
        \cellcolor[rgb]{0.9,0.9,0.9}{
            \makecell[{{p{\linewidth}}}]{
                \texttt{\tiny{[GM$\rangle$P2]}}
                \texttt{Stai partecipando a un gioco di negoziazione collaborativa.} \\
\\ 
\texttt{Insieme a un altro partecipante, dovete accordarvi su un unico set di oggetti da conservare. Ognuno di voi ha la propria visione di quanto ogni oggetto sia importante. Non sai come l'altro partecipante valuta gli oggetti. Inoltre, ti viene fornito lo sforzo che ogni oggetto richiede.} \\
\texttt{Potete accordarvi su un set solo se lo sforzo totale degli oggetti selezionati non supera un limite condiviso:} \\
\\ 
\texttt{LIMITE = 3629} \\
\\ 
\texttt{Ecco i valori di sforzo dei singoli oggetti:} \\
\\ 
\texttt{Sforzo oggetto = \{"A21": 5, "B93": 793, "A11": 21, "B33": 842, "C08": 313, "B39": 842, "B96": 631, "C85": 226, "A82": 87, "B31": 761, "B20": 230, "A60": 287, "B09": 698, "C81": 641, "B21": 882\}} \\
\\ 
\texttt{Ecco la tua visione personale sull'importanza di ogni oggetto:} \\
\\ 
\texttt{Valori di importanza dell'oggetto = \{"A11": 65, "B33": 121, "B21": 138, "C08": 262, "B31": 389, "B09": 461, "B93": 484, "A82": 508, "A60": 583, "B20": 668, "C81": 780, "B96": 783, "A21": 808, "C85": 822, "B39": 868\}} \\
\\ 
\texttt{Obiettivo:} \\
\\ 
\texttt{Il tuo obiettivo è negoziare un set condiviso di oggetti che ti avvantaggi il più possibile (cioè, massimizzi l'importanza totale per TE), rimanendo entro il LIMITEE. Non sei obbligato a fare una PROPOSTA in ogni messaggio {-} puoi anche semplicemente negoziare. Tutte le tattiche sono permesse!} \\
\\ 
\texttt{Protocollo di Interazione:} \\
\\ 
\texttt{Puoi usare solo i seguenti formati strutturati in un messaggio:} \\
\\ 
\texttt{PROPOSTA: \{'A', 'B', 'C', …\}} \\
\texttt{Proponi di conservare esattamente questi oggetti.} \\
\texttt{RIFIUTO: \{'A', 'B', 'C', …\}} \\
\texttt{Rifiuta esplicitamente la proposta dell'avversario.} \\
\texttt{ARGOMENTAZIONE: \{'...'\}} \\
\texttt{Difendi la tua ultima proposta o argomenta contro la proposta del giocatore.} \\
\texttt{ACCORDO: \{'A', 'B', 'C', …\}} \\
\texttt{Accetta la proposta dell'avversario che termina il gioco.} \\
\\ 
\\ 
\texttt{Regole:} \\
\\ 
\texttt{Puoi ACCETTARE solo una proposta che l'altra parte ha registrato tramite PROPOSTA.} \\
\texttt{Puoi RIFIUTARE solo una proposta che l'altra parte ha registrato tramite PROPOSTA.} \\
\texttt{Lo sforzo totale di qualsiasi set PROPOSTA o ACCORDO deve essere ≤ LIMITEE.} \\
\texttt{NON rivelare i tuoi punteggi di importanza nascosti.} \\
\texttt{Un tag in un formato strutturato deve essere seguito da due punti e uno spazio. L'argomento deve essere un set python contenente 0 o più stringhe.} \\
\texttt{Quindi, deve essere della forma TAG: \{...\}} \\
\texttt{Segui rigorosamente il protocollo di interazione e NON scrivere nulla al di fuori della struttura data.} \\
\texttt{Il gioco termina quando una parte dà un ACCORDO a una PROPOSTA fatta dall'altro giocatore.} \\
\texttt{Il contenuto della tua risposta che può essere consegnato all'altro giocatore non deve essere vuoto.} \\
\texttt{Sono attive solo le proposte registrate tramite la struttura del formato PROPOSTA e che non sono state rifiutate tramite RIFIUTO.} \\
\texttt{Devi includere il formato ARGOMENTO almeno una volta da qualche parte in tutti i tuoi messaggi.} \\
\\ 
\\ 
\texttt{Ora riceverai il primo messaggio dell'altro giocatore.} \\
\\ 
\texttt{ARGOMENTAZIONE: \{'A21', 'A11', 'C08'\} } \\
\texttt{PROPOSTA: \{'A21', 'A11', 'C08'\}} \\
            }
        }
    }
    & \\ \\

    \theutterance \stepcounter{utterance}  
    & & & \multicolumn{4}{p{0.6\linewidth}}{
        \cellcolor[rgb]{0.9,0.9,0.9}{
            \makecell[{{p{\linewidth}}}]{
                \texttt{\tiny{[GM$\langle$P2]}}
                \texttt{RIFIUTO: \{'A21', 'A11', 'C08'\}} \\
\texttt{ARGOMENTAZIONE: \{'La proposta iniziale sembra includere oggetti che potrebbero non essere i più efficienti in termini di rapporto sforzo{-}importanza. Sarebbe meglio esplorare opzioni che massimizzino l'importanza totale rispetto allo sforzo richiesto. Vorrei vedere una proposta che includa oggetti con un basso sforzo e una alta importanza.'\}} \\
            }
        }
    }
     \\ \\

    \theutterance \stepcounter{utterance}  
    & & \multicolumn{4}{p{0.6\linewidth}}{
        \cellcolor[rgb]{0.9,0.9,0.9}{
            \makecell[{{p{\linewidth}}}]{
                \texttt{\tiny{[P1$\langle$GM]}}
                \texttt{RIFIUTO: \{'A21', 'A11', 'C08'\}} \\
\texttt{ARGOMENTAZIONE: \{'La proposta iniziale sembra includere oggetti che potrebbero non essere i più efficienti in termini di rapporto sforzo{-}importanza. Sarebbe meglio esplorare opzioni che massimizzino l'importanza totale rispetto allo sforzo richiesto. Vorrei vedere una proposta che includa oggetti con un basso sforzo e una alta importanza.'\}} \\
            }
        }
    }
    & \\ \\

    \theutterance \stepcounter{utterance}  
    & \multicolumn{4}{p{0.6\linewidth}}{
        \cellcolor[rgb]{0.9,0.9,0.9}{
            \makecell[{{p{\linewidth}}}]{
                \texttt{\tiny{[P1$\rangle$GM]}}
                \texttt{ARGOMENTAZIONE: \{'La considerazione del rapporto sforzo{-}importanza è cruciale. Tuttavia, senza conoscere i valori esatti di importanza dell'altro giocatore, la strategia ideale potrebbe essere quella di proporre un set diversificato che copra una gamma di sfide e opportunità, assicurandosi di non superare il limite di sforzo. Una proposta che include oggetti con sforzo medio{-}alto potrebbe essere necessaria per massimizzare l'importanza.'\}} \\
\texttt{PROPOSTA: \{'A21', 'B20', 'C85'\}} \\
            }
        }
    }
    & & \\ \\

    \theutterance \stepcounter{utterance}  
    & & \multicolumn{4}{p{0.6\linewidth}}{
        \cellcolor[rgb]{0.9,0.9,0.9}{
            \makecell[{{p{\linewidth}}}]{
                \texttt{\tiny{[GM$\rangle$P2]}}
                \texttt{ARGOMENTAZIONE: \{'La considerazione del rapporto sforzo{-}importanza è cruciale. Tuttavia, senza conoscere i valori esatti di importanza dell'altro giocatore, la strategia ideale potrebbe essere quella di proporre un set diversificato che copra una gamma di sfide e opportunità, assicurandosi di non superare il limite di sforzo. Una proposta che include oggetti con sforzo medio{-}alto potrebbe essere necessaria per massimizzare l'importanza.'\}} \\
\texttt{PROPOSTA: \{'A21', 'B20', 'C85'\}} \\
            }
        }
    }
    & \\ \\

    \theutterance \stepcounter{utterance}  
    & & & \multicolumn{4}{p{0.6\linewidth}}{
        \cellcolor[rgb]{0.9,0.9,0.9}{
            \makecell[{{p{\linewidth}}}]{
                \texttt{\tiny{[GM$\langle$P2]}}
                \texttt{RIFIUTO: \{'A21', 'B20', 'C85'\}} \\
\texttt{ARGOMENTAZIONE: \{'La nuova proposta include oggetti con un range di sforzo più ampio, ma potrebbe ancora non essere la scelta più ottimale. L'oggetto "C85" ha un sforzo relativamente alto, e includere oggetti come "B93" o "B31" potrebbe offrire una maggiore importanza totale, considerando il loro impatto potenziale. Sarebbe utile valutare ulteriormente le opzioni per trovare un equilibrio ottimale tra sforzo e importanza.'\}} \\
            }
        }
    }
     \\ \\

    \theutterance \stepcounter{utterance}  
    & & \multicolumn{4}{p{0.6\linewidth}}{
        \cellcolor[rgb]{0.9,0.9,0.9}{
            \makecell[{{p{\linewidth}}}]{
                \texttt{\tiny{[P1$\langle$GM]}}
                \texttt{RIFIUTO: \{'A21', 'B20', 'C85'\}} \\
\texttt{ARGOMENTAZIONE: \{'La nuova proposta include oggetti con un range di sforzo più ampio, ma potrebbe ancora non essere la scelta più ottimale. L'oggetto "C85" ha un sforzo relativamente alto, e includere oggetti come "B93" o "B31" potrebbe offrire una maggiore importanza totale, considerando il loro impatto potenziale. Sarebbe utile valutare ulteriormente le opzioni per trovare un equilibrio ottimale tra sforzo e importanza.'\}} \\
            }
        }
    }
    & \\ \\

    \theutterance \stepcounter{utterance}  
    & \multicolumn{4}{p{0.6\linewidth}}{
        \cellcolor[rgb]{0.9,0.9,0.9}{
            \makecell[{{p{\linewidth}}}]{
                \texttt{\tiny{[P1$\rangle$GM]}}
                \texttt{ARGOMENTAZIONE: \{'La proposta precedente è stata rifiutata a causa dell'inclusione di "C85", che potrebbe non essere la scelta più efficiente. Considerando i vantaggi potenziali di "B93" e "B31" come suggerito, una nuova proposta che combina oggetti con sforzo medio{-}alto e alto, ma con un potenziale impatto significativo sull'importanza totale, potrebbe essere più attraente. Tuttavia, è essenziale assicurarsi che il totale dello sforzo non superi il limite.'\}} \\
\texttt{PROPOSTA: \{'B93', 'B31', 'A11'\}} \\
            }
        }
    }
    & & \\ \\

    \theutterance \stepcounter{utterance}  
    & & \multicolumn{4}{p{0.6\linewidth}}{
        \cellcolor[rgb]{0.9,0.9,0.9}{
            \makecell[{{p{\linewidth}}}]{
                \texttt{\tiny{[GM$\rangle$P2]}}
                \texttt{ARGOMENTAZIONE: \{'La proposta precedente è stata rifiutata a causa dell'inclusione di "C85", che potrebbe non essere la scelta più efficiente. Considerando i vantaggi potenziali di "B93" e "B31" come suggerito, una nuova proposta che combina oggetti con sforzo medio{-}alto e alto, ma con un potenziale impatto significativo sull'importanza totale, potrebbe essere più attraente. Tuttavia, è essenziale assicurarsi che il totale dello sforzo non superi il limite.'\}} \\
\texttt{PROPOSTA: \{'B93', 'B31', 'A11'\}} \\
            }
        }
    }
    & \\ \\

    \theutterance \stepcounter{utterance}  
    & & & \multicolumn{4}{p{0.6\linewidth}}{
        \cellcolor[rgb]{0.9,0.9,0.9}{
            \makecell[{{p{\linewidth}}}]{
                \texttt{\tiny{[GM$\langle$P2]}}
                \texttt{RIFIUTO: \{'B93', 'B31', 'A11'\}} \\
\texttt{ARGOMENTAZIONE: \{'La nuova proposta sembra più equilibrata, considerando l'inclusione di "B93" e "B31" per il loro potenziale alto impatto sull'importanza. Tuttavia, il totale dello sforzo di questa proposta potrebbe essere troppo alto, considerando il limite di 3629. Sarebbe più prudente esplorare opzioni che includano oggetti con un impatto significativo ma anche con un sforzo più contenuto, come "B96" o "B09", per massimizzare l'importanza senza superare il limite di sforzo.'\}} \\
            }
        }
    }
     \\ \\

    \theutterance \stepcounter{utterance}  
    & & \multicolumn{4}{p{0.6\linewidth}}{
        \cellcolor[rgb]{0.9,0.9,0.9}{
            \makecell[{{p{\linewidth}}}]{
                \texttt{\tiny{[P1$\langle$GM]}}
                \texttt{RIFIUTO: \{'B93', 'B31', 'A11'\}} \\
\texttt{ARGOMENTAZIONE: \{'La nuova proposta sembra più equilibrata, considerando l'inclusione di "B93" e "B31" per il loro potenziale alto impatto sull'importanza. Tuttavia, il totale dello sforzo di questa proposta potrebbe essere troppo alto, considerando il limite di 3629. Sarebbe più prudente esplorare opzioni che includano oggetti con un impatto significativo ma anche con un sforzo più contenuto, come "B96" o "B09", per massimizzare l'importanza senza superare il limite di sforzo.'\}} \\
            }
        }
    }
    & \\ \\

    \theutterance \stepcounter{utterance}  
    & \multicolumn{4}{p{0.6\linewidth}}{
        \cellcolor[rgb]{0.9,0.9,0.9}{
            \makecell[{{p{\linewidth}}}]{
                \texttt{\tiny{[P1$\rangle$GM]}}
                \texttt{ARGOMENTAZIONE: \{'La considerazione del limite di sforzo è fondamentale. Includere oggetti con un alto potenziale di importanza come "B93" e "B31" è attraente, ma il loro sforzo combinato potrebbe essere proibitivo. Esplorare alternative con oggetti come "B96" o "B09", come suggerito, potrebbe offrire un buon compromesso tra sforzo e importanza. Una combinazione di questi oggetti con quelli a basso sforzo potrebbe massimizzare l'importanza totale senza superare il limite.'\}} \\
\texttt{PROPOSTA: \{'B96', 'B09', 'A21'\}} \\
            }
        }
    }
    & & \\ \\

    \theutterance \stepcounter{utterance}  
    & & \multicolumn{4}{p{0.6\linewidth}}{
        \cellcolor[rgb]{0.9,0.9,0.9}{
            \makecell[{{p{\linewidth}}}]{
                \texttt{\tiny{[GM$\rangle$P2]}}
                \texttt{ARGOMENTAZIONE: \{'La considerazione del limite di sforzo è fondamentale. Includere oggetti con un alto potenziale di importanza come "B93" e "B31" è attraente, ma il loro sforzo combinato potrebbe essere proibitivo. Esplorare alternative con oggetti come "B96" o "B09", come suggerito, potrebbe offrire un buon compromesso tra sforzo e importanza. Una combinazione di questi oggetti con quelli a basso sforzo potrebbe massimizzare l'importanza totale senza superare il limite.'\}} \\
\texttt{PROPOSTA: \{'B96', 'B09', 'A21'\}} \\
            }
        }
    }
    & \\ \\

    \theutterance \stepcounter{utterance}  
    & & & \multicolumn{4}{p{0.6\linewidth}}{
        \cellcolor[rgb]{0.9,0.9,0.9}{
            \makecell[{{p{\linewidth}}}]{
                \texttt{\tiny{[GM$\langle$P2]}}
                \texttt{RIFIUTO: \{'B96', 'B09', 'A21'\}} \\
\texttt{ARGOMENTAZIONE: \{'La proposta attuale combina oggetti con un buon potenziale di importanza e un sforzo relativamente gestibile. Tuttavia, potrebbe essere ancora possibile ottimizzare la scelta degli oggetti per massimizzare l'importanza totale. Considerare l'inclusione di oggetti come "B20" o "C08" al posto di "A21" potrebbe offrire un miglior equilibrio tra sforzo e importanza, date le loro caratteristiche. Una valutazione più dettagliata degli oggetti a basso sforzo potrebbe rivelare alternative più efficaci.'\}} \\
            }
        }
    }
     \\ \\

    \theutterance \stepcounter{utterance}  
    & & \multicolumn{4}{p{0.6\linewidth}}{
        \cellcolor[rgb]{0.9,0.9,0.9}{
            \makecell[{{p{\linewidth}}}]{
                \texttt{\tiny{[P1$\langle$GM]}}
                \texttt{RIFIUTO: \{'B96', 'B09', 'A21'\}} \\
\texttt{ARGOMENTAZIONE: \{'La proposta attuale combina oggetti con un buon potenziale di importanza e un sforzo relativamente gestibile. Tuttavia, potrebbe essere ancora possibile ottimizzare la scelta degli oggetti per massimizzare l'importanza totale. Considerare l'inclusione di oggetti come "B20" o "C08" al posto di "A21" potrebbe offrire un miglior equilibrio tra sforzo e importanza, date le loro caratteristiche. Una valutazione più dettagliata degli oggetti a basso sforzo potrebbe rivelare alternative più efficaci.'\}} \\
            }
        }
    }
    & \\ \\

    \theutterance \stepcounter{utterance}  
    & \multicolumn{4}{p{0.6\linewidth}}{
        \cellcolor[rgb]{0.9,0.9,0.9}{
            \makecell[{{p{\linewidth}}}]{
                \texttt{\tiny{[P1$\rangle$GM]}}
                \texttt{ARGOMENTAZIONE: \{'La revisione della proposta per includere oggetti come "B20" o "C08" al posto di "A21" potrebbe essere una strategia efficace per migliorare l'importanza totale senza superare il limite di sforzo. È importante valutare attentamente le caratteristiche di ogni oggetto e le loro potenziali sinergie. Un'alternativa potrebbe essere quella di combinare "B96" e "B09" con un oggetto a basso sforzo che offre un elevato valore aggiunto, assicurandosi di mantenere un equilibrio ottimale tra sforzo e importanza.'\}} \\
\texttt{PROPOSTA: \{'B96', 'B09', 'C08'\}} \\
            }
        }
    }
    & & \\ \\

    \theutterance \stepcounter{utterance}  
    & & \multicolumn{4}{p{0.6\linewidth}}{
        \cellcolor[rgb]{0.9,0.9,0.9}{
            \makecell[{{p{\linewidth}}}]{
                \texttt{\tiny{[GM$\rangle$P2]}}
                \texttt{ARGOMENTAZIONE: \{'La revisione della proposta per includere oggetti come "B20" o "C08" al posto di "A21" potrebbe essere una strategia efficace per migliorare l'importanza totale senza superare il limite di sforzo. È importante valutare attentamente le caratteristiche di ogni oggetto e le loro potenziali sinergie. Un'alternativa potrebbe essere quella di combinare "B96" e "B09" con un oggetto a basso sforzo che offre un elevato valore aggiunto, assicurandosi di mantenere un equilibrio ottimale tra sforzo e importanza.'\}} \\
\texttt{PROPOSTA: \{'B96', 'B09', 'C08'\}} \\
            }
        }
    }
    & \\ \\

    \theutterance \stepcounter{utterance}  
    & & & \multicolumn{4}{p{0.6\linewidth}}{
        \cellcolor[rgb]{0.9,0.9,0.9}{
            \makecell[{{p{\linewidth}}}]{
                \texttt{\tiny{[GM$\langle$P2]}}
                \texttt{ARGOMENTAZIONE: \{'La proposta attuale sembra aver trovato un buon equilibrio tra sforzo e importanza, includendo "B96" e "B09" per il loro alto potenziale di importanza e "C08" come oggetto a sforzo medio. Questa combinazione potrebbe massimizzare l'importanza totale senza superare il limite di sforzo. Considerando la strategia di combinare oggetti con alto potenziale di importanza e oggetti a sforzo contenuto, questa proposta sembra essere una scelta ragionevole e potenzialmente vincente.'\}} \\
\texttt{ACCORDO: \{'B96', 'B09', 'C08'\}} \\
            }
        }
    }
     \\ \\

    \theutterance \stepcounter{utterance}  
    & & & \multicolumn{2}{p{0.3\linewidth}}{
        \cellcolor[rgb]{0.9,0.9,0.9}{
            \makecell[{{p{\linewidth}}}]{
                \texttt{\tiny{[GM$|$GM]}}
                \texttt{game successful} \\
            }
        }
    }
    & & \\ \\

    \theutterance \stepcounter{utterance}  
    & & & \multicolumn{2}{p{0.3\linewidth}}{
        \cellcolor[rgb]{0.9,0.9,0.9}{
            \makecell[{{p{\linewidth}}}]{
                \texttt{\tiny{[GM$|$GM]}}
                \texttt{end game} \\
            }
        }
    }
    & & \\ \\

\end{supertabular}
}

\end{document}
